% Modelo de relatório de experimento em formato de artigo científico

% Criado a partir de "Modelo de Relatório abnTeX", de autoria de Gustavo Yudi Bientinezi Matsuzake, disponibilizado na Galeria Overleaf (https://pt.overleaf.com/latex/templates/modelo-de-relatorio-abntex/cffmvqrhsprn).

% Definição do documento e pacotes utilizados. 

\documentclass[12pt,a4paper,twoside]{article}

% Configurações:

\input{conf/texto.tex}
\input{conf/matematica.tex}
\input{conf/ciencia.tex}
\input{conf/ambientes.tex}

% Para citações:

\usepackage[brazilian,hyperpageref]{backref}
\usepackage[alf]{abntex2cite}

% EDITE A PARTIR DAQUI:

\def \tituloportugues {COLOQUE SEU TÍTULO EM PORTUGUÊS AQUI, ELE PASSARÁ PARA A OUTRA LINHA AUTOMATICAMENTE SE PRECISO}  % Título. [EDITAR]
\def \tituloingles {COLOQUE O TÍTULO EM INGLÊS AQUI, ELE PASSARÁ PARA A OUTRA LINHA AUTOMATICAMENTE SE PRECISO}  % Título em inglês. [EDITAR]
\setcounter{page}{1} % página inicial do relatório
\def \tipoartigo {Grupo 1}
\def \tema {Física Moderna} % pode mudar para "Eletromagnetismo", "Física Quântica" etc.
\def \coordenadas {\textit{Laboratório de Física 4 experimental do IFUnB}, Brasília, 2019.} % Título em inglês. [EDITAR]
\def \paginainternet {https://ifserv.fis.unb.br/moodle}
\definecolor{cortextoquadro}{RGB}{0,132,61} % Cor do texto dos quadros do leiaute.
\definecolor{corfundoquadro}{RGB}{200,200,200} % Cor de fundo dos quadros do leiaute.
\def \fontetexto {\bfseries} % Configurações de fonte para textos dentro dos quadros do leiaute.
\def \logorevista {figs/logo-unb.png} % Defina o arquivo do logotipo.
\def \larguralogo {2.3cm} % Defina a largura com que o logotipo aparecerá na capa.

% As configurações de layout estão no arquivo mostrado abaixo:
\input{conf/layout.tex} % Não mexa nesta linha!

% Coloque abaixo os nomes dos autores e o número de suas instituições:

\author[1]{Coloque Nome do Autor Um}
\author[2]{Coloque Nome do Autor Dois}
\author[1]{Coloque Nome do Autor Três}
\author[2]{Coloque Nome do Autor Quatro}

% Coloque abaixo os números e nomes das instituições:

\affil[1]{Coloque aqui o nome da instituição 1.}
\affil[2]{Coloque aqui o nome da instituição 2.}
\affil[3]{Coloque aqui o nome da instituição 3.}

% Ajustes diversos antes de começar o documento (últimos preparativos),
% tais como atribuição da variável "title", configurações para resolver erros de compilação etc.

\input{conf/ultimos-preparativos.tex}

%-----------------------------

% Aqui damos início ao documento.

\begin{document}

\maketitle

\thispagestyle{capa}

% Resumo e palavras-chave. 

\hrule

\selectlanguage{portuges}
\begin{abstract}
\noindent
\input{texto/resumo.tex} % [EDITAR] %
\vspace{\baselineskip}
\palavraschave{física, professor, didática.} % [EDITAR] %
\end{abstract}
\vspace{5mm}

\hrule

\selectlanguage{english}
\begin{abstract}
\noindent
\input{texto/abstract.tex} % [EDITAR] %
\vspace{\baselineskip}
\keywords{physics, teacher, didactic method.} % [EDITAR] %
\end{abstract}
\vspace{5mm}

\hrule

\selectlanguage{portuges}

\section{Instruções}

    \input{texto/instrucoes.tex}

\section{Introdução} % [EDITAR]

    \input{texto/introducao.tex} % [EDITAR]

\section{Objetivo} % [EDITAR]

    \input{texto/objetivos.tex} % [EDITAR]

\section{Materiais} % [EDITAR]

    \input{texto/materiais.tex} % [EDITAR]

\section{Procedimento} % [EDITAR]

    \input{texto/procedimento.tex} % [EDITAR]

\section{Resultados} % [EDITAR]

\subsection{Dados}

    \input{texto/dados.tex} % [EDITAR]

\subsection{Análise} % [EDITAR]

    \input{texto/analise.tex} % [EDITAR]

\section{Discussao} % [EDITAR]

    \input{texto/discussao.tex} % [EDITAR]

\section{Conclusão} % [EDITAR]

    \input{texto/conclusao.tex} % [EDITAR]


\section*{Referências} % [EDITAR]

    \noindent Paardekooper S.-J. et al. (2010) Mon. Not. R. Astron. Soc., 401,
1950. % [EDITAR]

    \noindent Duncan, M. J., Levison, H. F. \& Lee, M.H. (1998). A.J., 116, 2067. % [EDITAR]

    \noindent Baruteau et al. (2014) in: Protoplanetary Disks and Stars VI, University of Arizona Press, 667. % [EDITAR]

\vspace{\baselineskip} \hrule

\vspace{2\baselineskip}

Se você utilizar a seção acima, apague o comando ``bibliography{main}'' do código. A referência também pode ser criada usando Bibitem, BibTeX, BibLaTeX etc. Caso você use BibTeX, apague a seção acima, insira as referências no arquivo ``main.bib''. Depois cite segundo este modelo: \cite{greenwade93}. Ou, para citar várias referências juntas, assim: \cite{greenwade93,wikimedia-latex}. Suas referências ficarão assim:

\bibliography{relatorio}

\end{document}

% Tutoriais LaTeX:

% <http://www.mat.ufmg.br/~regi/topicos/intlat.pdf>
% <http://ece.uprm.edu/~caceros/latex/introduction.pdf>

%-----------------------------

% LaTeX Online (free account):

% <https://www.overleaf.com/>

%-----------------------------

% Códigos matemáticos LaTeX:

% <http://www.codecogs.com/latex/eqneditor.php>

%-----------------------------

% Construir tabelas no LaTeX:

% <http://www.tablesgenerator.com/#>

%-----------------------------